\section{Bab 14: Skalabilitas, Sistem Web Terdistribusi, dan Emerging Web}

\subsection*{Tujuan Pembelajaran}
\begin{itemize}
  \item Mengukur penskalaan horizontal sampai satu komponen bersama menjadi
        \textit{bottleneck}, lalu menunjukkan pergeseran
        \textit{bottleneck}-nya dari angka hasil pengukuran.
  \item Merancang rencana kapasitas untuk beban 100 kali lipat yang mencakup
        batas layanan, replikasi, partisi, dan biaya.
  \item Membandingkan manfaat dan biaya tiga teknologi yang sedang tumbuh,
        yaitu mode \textit{offline} Progressive Web App (PWA), \textit{cold
        start} pada \textit{serverless}, dan fitur berbantuan Large Language
        Model (LLM), berdasarkan hasil pengukuran.
\end{itemize}

Seluruh file yang dipakai pada latihan ini tersedia di
\berkas{sources/chapter14/}, dihitung relatif terhadap akar repositori.
Latihan ini hanya membutuhkan Python 3 dan Google Chrome atau Chromium, tanpa
Docker dan tanpa pustaka tambahan. Skrip pengukur menyalakan dan mematikan
sendiri proses yang diukurnya.

\berkas{ukur-skala.py} sengaja membebani, tetapi yang dibebaninya hanya
proses yang dinyalakannya sendiri di 127.0.0.1. Seluruh alat ukur lain
bersifat mengamati. Angka \textit{throughput} bergantung pada jumlah \textit{core} dan
batas daya komputer, sehingga hasil tiap mahasiswa akan berbeda,
sedangkan pola perbandingannya akan sama. Angka pada bab ini diukur dengan
Python 3.15 dan Chrome 157 pada laptop Intel Core i9-11900H.

Sebelum mulai, periksa kedua perkakasnya seperti pada Listing
\ref{lst:c14-siapkan}. \textit{Sanity check} ini mencegah pengukuran yang
gagal di tengah jalan karena Chrome tidak ditemukan.

\begin{lstlisting}[language=bash, caption={Memeriksa perkakas latihan, dijalankan dari \berkas{sources/chapter14/}}, label={lst:c14-siapkan}]
cd sources/chapter14
python3 --version
google-chrome --version

# Keluarannya:
# Python 3.15.2
# Google Chrome 157.0.8103.38
\end{lstlisting}

\subsection*{Latihan 1: Penskalaan Horizontal dan Bottleneck yang Pindah}

Latihan ini membuktikan bahwa menambah replika berhenti menolong ketika
komponen bersamanya penuh, lalu memperlihatkan \textit{bottleneck} pindah
ketika komponen itu dibesarkan.

\begin{enumerate}
  \item Jalankan \berkas{ukur-skala.py}, seperti pada Listing
        \ref{lst:c14-lat1-jalan}, lalu catat keempat barisnya.
  \item Kalikan \textit{throughput} dengan p50 dalam detik pada tiap baris,
        lalu bandingkan hasilnya dengan jumlah \textit{worker}, yaitu 32.
  \item Jalankan ulang dengan \texttt{JUMLAH\_KONEKSI=16} di depan
        perintahnya, lalu catat keempat barisnya.
  \item Isi Tabel \ref{tab:c14-lat1-hasil}. Tabel \ref{tab:c14-lat1-contoh}
        memperlihatkan contoh yang sudah terisi.
  \item Jawab tiga pertanyaan berikut. Pertama, mengapa empat dan delapan
        replika berhenti pada angka yang hampir sama dengan empat koneksi?
        Kedua, mengapa p50 empat dan delapan replika sama persis, padahal
        replikanya berlipat? Ketiga, mengapa delapan replika dengan 16
        koneksi hanya mencapai efisiensi sekitar 40\%, dan apa yang berubah
        bila tiap replika berjalan di mesin terpisah?
\end{enumerate}

\begin{lstlisting}[language=bash, caption={Mengukur penskalaan dengan dua kapasitas basis data, tersedia di \berkas{sources/chapter14/ukur-skala.py}}, label={lst:c14-lat1-jalan}]
python3 ukur-skala.py
JUMLAH_KONEKSI=16 python3 ukur-skala.py

# Keluarannya, baris tabel saja, 4 koneksi lalu 16 koneksi:
#   Replika  Permintaan/detik   p50 ms   p95 ms  Efisiensi
#         1             106.8    286.1    547.5       100%
#         2             186.6    161.4    278.3        87%
#         4             198.2    161.4    168.6        46%
#         8             198.0    161.4    168.9        23%
#
#         1             105.0    288.9    581.5       100%
#         2             179.8    169.5    302.2        86%
#         4             285.0    105.2    181.5        68%
#         8             347.8     87.4    143.3        41%
\end{lstlisting}

\begin{table}[!htb]
  \centering
  \caption{Contoh hasil Latihan 1 yang sudah terisi}
  \label{tab:c14-lat1-contoh}
  \begin{tabularx}{\textwidth}{Xlll}
    \hline
    \textbf{Replika} & \textbf{4 koneksi, per detik} & \textbf{16 koneksi, per detik} & \textbf{16 koneksi, efisiensi} \\
    \hline
    1 & 106,8 & 105,0 & 100\% \\
    \hline
    2 & 186,6 & 179,8 & 86\% \\
    \hline
    4 & 198,2 & 285,0 & 68\% \\
    \hline
    8 & 198,0 & 347,8 & 41\% \\
    \hline
  \end{tabularx}
\end{table}

\begin{table}[!htb]
  \centering
  \caption{Lembar isian hasil pengukuran Latihan 1}
  \label{tab:c14-lat1-hasil}
  \begin{tabularx}{\textwidth}{Xlll}
    \hline
    \textbf{Replika} & \textbf{4 koneksi, per detik} & \textbf{16 koneksi, per detik} & \textbf{16 koneksi, efisiensi} \\
    \hline
    1 & \dots & \dots & \dots \\
    \hline
    2 & \dots & \dots & \dots \\
    \hline
    4 & \dots & \dots & \dots \\
    \hline
    8 & \dots & \dots & \dots \\
    \hline
  \end{tabularx}
\end{table}

\subsection*{Latihan 2: Partisi Data dan Kunci yang Pindah}

Latihan ini membuktikan bahwa cara membagi data menentukan berapa banyak data
yang harus disalin ketika \textit{node} ditambah, dan bahwa angka yang kecil belum
tentu baik.

\begin{enumerate}
  \item Jalankan \berkas{partisi.py}, seperti pada Listing
        \ref{lst:c14-lat2-jalan}, lalu catat ketiga barisnya.
  \item Ubah \texttt{JUMLAH\_NODE\_AKHIR} di \berkas{partisi.py} menjadi 8,
        jalankan ulang, lalu catat kolom kunci pindahnya.
  \item Kembalikan \texttt{JUMLAH\_NODE\_AKHIR} menjadi 5.
  \item Isi Tabel \ref{tab:c14-lat2-hasil}. Tabel \ref{tab:c14-lat2-contoh}
        memperlihatkan contoh yang sudah terisi.
  \item Jawab tiga pertanyaan berikut. Pertama, mengapa modulo dari empat ke
        delapan \textit{node} hanya memindahkan 50,0\% kunci, jauh lebih sedikit
        daripada dari empat ke lima? Kedua, mengapa satu titik per \textit{node} yang
        hanya memindahkan 2,1\% kunci justru tanda buruk? Ketiga, mengapa
        \textit{node} terberat menurut permintaan selalu lebih berat daripada menurut
        kunci, berapa pun cara baginya?
\end{enumerate}

\begin{lstlisting}[language=bash, caption={Membandingkan dua cara membagi data, tersedia di \berkas{sources/chapter14/partisi.py}}, label={lst:c14-lat2-jalan}]
python3 partisi.py

# Keluarannya:
#   100000 kunci, 4 node menjadi 5 node, ideal 20.0% tiap node
#
#   Cara                           Pindah Terberat, kunci Terberat, permintaan
#   modulo                          80.2%           20.2%                24.3%
#   consistent hashing, 1 titik      2.1%           49.2%                43.3%
#   consistent hashing, 100 titik   20.2%           21.2%                26.9%
#
# Dengan JUMLAH_NODE_AKHIR = 8, tiga baris tabelnya:
#   modulo                          50.0%           12.6%                20.9%
#   consistent hashing, 1 titik     43.4%           27.7%                24.0%
#   consistent hashing, 100 titik   49.5%           13.7%                20.5%
\end{lstlisting}

\begin{table}[!htb]
  \centering
  \caption{Contoh hasil Latihan 2 yang sudah terisi}
  \label{tab:c14-lat2-contoh}
  \begin{tabularx}{\textwidth}{Xlll}
    \hline
    \textbf{Cara} & \textbf{Pindah, 4 ke 5} & \textbf{Pindah, 4 ke 8} & \textbf{Terberat, 4 ke 5} \\
    \hline
    modulo & 80,2\% & 50,0\% & 20,2\% \\
    \hline
    \textit{consistent hashing}, 1 titik & 2,1\% & 43,4\% & 49,2\% \\
    \hline
    \textit{consistent hashing}, 100 titik & 20,2\% & 49,5\% & 21,2\% \\
    \hline
  \end{tabularx}
\end{table}

\begin{table}[!htb]
  \centering
  \caption{Lembar isian hasil perhitungan Latihan 2}
  \label{tab:c14-lat2-hasil}
  \begin{tabularx}{\textwidth}{Xlll}
    \hline
    \textbf{Cara} & \textbf{Pindah, 4 ke 5} & \textbf{Pindah, 4 ke 8} & \textbf{Terberat, 4 ke 5} \\
    \hline
    modulo & \dots & \dots & \dots \\
    \hline
    \textit{consistent hashing}, 1 titik & \dots & \dots & \dots \\
    \hline
    \textit{consistent hashing}, 100 titik & \dots & \dots & \dots \\
    \hline
  \end{tabularx}
\end{table}

Angka pada latihan ini dihitung, bukan diukur dengan jam, sehingga akan sama
persis pada komputer mana pun. Yang dilaporkan tetap pola perbandingannya,
karena angka itulah yang menjadi dasar memilih cara bagi.

\subsection*{Latihan 3: Rencana Kapasitas untuk Beban 100 Kali Lipat}

Latihan ini menyusun rencana penskalaan dari angka hasil pengukuran sendiri,
lalu memperlihatkan komponen mana yang paling mahal dan mana yang pertama kali
melewati batas.

\begin{enumerate}
  \item Isi \texttt{kapasitas\_replika\_rps} di \berkas{rencana.json} dengan
        \textit{throughput} satu replika dari Latihan 1, dan
        \texttt{kapasitas\_node\_basis\_data\_qps} dengan
        \textit{throughput} empat replika pada empat koneksi.
  \item Jalankan \berkas{rencana-kapasitas.py}, seperti pada Listing
        \ref{lst:c14-lat3-jalan}, lalu catat baris totalnya.
  \item Ubah \texttt{hit\_ratio\_cache} katalog menjadi 0,0, jalankan ulang,
        catat totalnya, lalu kembalikan menjadi 0,9.
  \item Ubah \texttt{kelipatan\_beban} menjadi 200, jalankan ulang, catat
        pemeriksaan koneksi yang gagal, lalu kembalikan menjadi 100.
  \item Isi Tabel \ref{tab:c14-lat3-hasil}. Tabel \ref{tab:c14-lat3-contoh}
        memperlihatkan contoh yang sudah terisi.
  \item Jawab tiga pertanyaan berikut. Pertama, mengapa pesanan membutuhkan
        lima \textit{shard} sedangkan katalog cukup satu, padahal beban
        katalog empat kali lebih besar? Kedua, berapa \textit{node} basis data dan
        berapa rupiah per bulan yang dihemat \textit{cache} katalog? Ketiga,
        pada beban 200 kali lipat, batas apa yang pertama kali terlewati, dan
        komponen apa yang menjawabnya?
\end{enumerate}

\begin{lstlisting}[language=bash, caption={Menyusun rencana kapasitas, tersedia di \berkas{sources/chapter14/rencana-kapasitas.py}}, label={lst:c14-lat3-jalan}]
python3 rencana-kapasitas.py

# Keluarannya:
#   Beban 100 kali lipat, utilisasi maksimum 60%
#
#   Layanan   Target Replika  Baca/s Tulis/s Shard Node BD Juta Rp
#   katalog     1200      19   235.2    48.0     1       3     9.3
#   pesanan      300       5   360.0   540.0     5      10    13.5
#   akun         200       4    90.0    20.0     1       2     3.6
#
#   katalog     76 koneksi per node basis data, batas 100: aman
#   pesanan     20 koneksi per node basis data, batas 100: aman
#   akun        16 koneksi per node basis data, batas 100: aman
#
#   Total: 28 replika, 15 node basis data, Rp 26.4 juta per bulan
\end{lstlisting}

\begin{table}[!htb]
  \centering
  \caption{Contoh hasil Latihan 3 yang sudah terisi}
  \label{tab:c14-lat3-contoh}
  \begin{tabularx}{\textwidth}{Xllll}
    \hline
    \textbf{Skenario} & \textbf{Replika} & \textbf{\textit{Node} basis data} & \textbf{Juta Rp} & \textbf{Koneksi katalog} \\
    \hline
    100 kali, \textit{cache} 90\% & 28 & 15 & 26,4 & 76, aman \\
    \hline
    100 kali, tanpa \textit{cache} katalog & 28 & 33 & 48,0 & 76, aman \\
    \hline
    200 kali, \textit{cache} 90\% & 55 & 28 & 50,1 & 152, perlu \textit{pooler} \\
    \hline
  \end{tabularx}
\end{table}

\begin{table}[!htb]
  \centering
  \caption{Lembar isian hasil perhitungan Latihan 3}
  \label{tab:c14-lat3-hasil}
  \begin{tabularx}{\textwidth}{Xllll}
    \hline
    \textbf{Skenario} & \textbf{Replika} & \textbf{\textit{Node} basis data} & \textbf{Juta Rp} & \textbf{Koneksi katalog} \\
    \hline
    100 kali, \textit{cache} 90\% & \dots & \dots & \dots & \dots \\
    \hline
    100 kali, tanpa \textit{cache} katalog & \dots & \dots & \dots & \dots \\
    \hline
    200 kali, \textit{cache} 90\% & \dots & \dots & \dots & \dots \\
    \hline
  \end{tabularx}
\end{table}

Biaya pada rencana ini hanyalah asumsi contoh yang tertulis di
\berkas{rencana.json}. Yang dinilai adalah susunan dan alasannya: layanan
mana yang membutuhkan \textit{shard}, berapa yang dihemat \textit{cache}, dan
batas mana yang terlewati lebih dahulu.

\subsection*{Latihan 4: Mode Offline Progressive Web App}

Latihan ini membuktikan bahwa halaman yang dilindungi \textit{service worker}
tetap tampil ketika servernya \textit{down}, dan bahwa perlindungan itu hanya berlaku
bagi pengguna yang pernah datang.

\begin{enumerate}
  \item Jalankan \berkas{ukur-offline.py}, seperti pada Listing
        \ref{lst:c14-lat4-jalan}, lalu catat keempat barisnya.
  \item Cocokkan jumlah permintaan kunjungan pertama dengan isi
        \texttt{KERANGKA} di \berkas{statis/sw.js}.
  \item Nyalakan \texttt{python3 aplikasi.py}, buka
        \texttt{http://127.0.0.1:8120/} di Chrome dua kali, lalu periksa
        bagian Application di DevTools: \textit{service worker}-nya dan isi
        \textit{cache} \texttt{katalog-v1}.
  \item Matikan \berkas{aplikasi.py}, muat ulang halamannya, lalu catat jam
        pembuatan data yang tertulis di atas daftar buku.
  \item Isi Tabel \ref{tab:c14-lat4-hasil}. Tabel \ref{tab:c14-lat4-contoh}
        memperlihatkan contoh yang sudah terisi.
  \item Jawab tiga pertanyaan berikut. Pertama, mengapa kunjungan pertama
        justru yang paling banyak mengirim permintaan? Kedua, kunjungan
        ketiga dan keempat sama-sama tanpa server, apa yang membuat hasilnya
        berbeda? Ketiga, apa yang wajib diperlihatkan halaman kepada pengguna
        ketika datanya berasal dari \textit{cache}?
\end{enumerate}

\begin{lstlisting}[language=bash, caption={Empat kunjungan dengan Chrome tanpa jendela, tersedia di \berkas{sources/chapter14/ukur-offline.py}}, label={lst:c14-lat4-jalan}]
python3 ukur-offline.py

# Keluarannya:
#   Kunjungan       Server Buku  Sumber    Saat mulai    Tampil ms Ke server
#   1, profil baru  hidup     8  jaringan  belum               377        11
#   2, profil sama  hidup     8  jaringan  mengendalikan        62         2
#   3, profil sama  down      8  cache     mengendalikan        31         0
#   4, profil baru  down      0  gagal     -                     -         0
#                   Chrome: net::ERR_CONNECTION_REFUSED
\end{lstlisting}

\begin{table}[!htb]
  \centering
  \caption{Contoh hasil Latihan 4 yang sudah terisi}
  \label{tab:c14-lat4-contoh}
  \begin{tabularx}{\textwidth}{Xllll}
    \hline
    \textbf{Kunjungan} & \textbf{Server} & \textbf{Buku} & \textbf{Sumber} & \textbf{Ke server} \\
    \hline
    1, profil baru & hidup & 8 & jaringan & 11 \\
    \hline
    2, profil sama & hidup & 8 & jaringan & 2 \\
    \hline
    3, profil sama & \textit{down} & 8 & \textit{cache} & 0 \\
    \hline
    4, profil baru & \textit{down} & 0 & gagal & 0 \\
    \hline
  \end{tabularx}
\end{table}

\begin{table}[!htb]
  \centering
  \caption{Lembar isian hasil pengamatan Latihan 4}
  \label{tab:c14-lat4-hasil}
  \begin{tabularx}{\textwidth}{Xllll}
    \hline
    \textbf{Kunjungan} & \textbf{Server} & \textbf{Buku} & \textbf{Sumber} & \textbf{Ke server} \\
    \hline
    1, profil baru & hidup & \dots & \dots & \dots \\
    \hline
    2, profil sama & hidup & \dots & \dots & \dots \\
    \hline
    3, profil sama & \textit{down} & \dots & \dots & \dots \\
    \hline
    4, profil baru & \textit{down} & \dots & \dots & \dots \\
    \hline
  \end{tabularx}
\end{table}

Jumlah permintaan pada kunjungan kedua bergantung pada versi Chrome. Chrome
157 yang dipakai pada bab ini mengirim satu permintaan navigasi bersamaan
dengan menyalakan \textit{service worker}, sehingga hasilnya dua, bukan satu.
Browser lain atau versi lain dapat mencatat angka berbeda, sedangkan pola
kunjungan ketiga dan keempat akan sama.

\subsection*{Latihan 5: Cold Start Fungsi Serverless}

Latihan ini membuktikan bahwa \textit{cold start} ditentukan oleh apa yang
dimuat sebelum fungsinya dipanggil, bukan oleh fungsinya sendiri.

\begin{enumerate}
  \item Jalankan \berkas{ukur-cold-start.py}, seperti pada Listing
        \ref{lst:c14-lat5-jalan}, lalu catat ketiga barisnya.
  \item Tambahkan \texttt{"multiprocessing"}, \texttt{"ssl"}, dan
        \texttt{"difflib"} ke \texttt{DAFTAR\_DEPENDENSI} di
        \berkas{fungsi.py}, jalankan ulang, lalu catat baris ketiganya.
  \item Kembalikan \texttt{DAFTAR\_DEPENDENSI} seperti semula.
  \item Isi Tabel \ref{tab:c14-lat5-hasil}. Tabel \ref{tab:c14-lat5-contoh}
        memperlihatkan contoh yang sudah terisi.
  \item Jawab tiga pertanyaan berikut. Pertama, mengapa \textit{cold start}
        lebih dari seribu kali panggilan \textit{warm}, padahal fungsinya
        sama? Kedua, bagian mana dari \textit{cold start} layanan
        \textit{serverless} sungguhan yang tidak terukur di sini? Ketiga,
        mengapa pekerjaan persiapan seperti membuka koneksi basis data
        sebaiknya ditaruh di luar \textit{handler}?
\end{enumerate}

\begin{lstlisting}[language=bash, caption={Membandingkan \textit{warm} dan \textit{cold start}, tersedia di \berkas{sources/chapter14/ukur-cold-start.py}}, label={lst:c14-lat5-jalan}]
python3 ukur-cold-start.py

# Keluarannya:
#   Cara                             p50 ms   p95 ms  Kali lipat
#   warm, proses dipakai ulang        0.012    0.012           1
#   cold start                       19.069   22.248        1610
#   cold start, dengan dependensi   105.565  118.585        8911
#
# Baris ketiga dengan tiga dependensi tambahan:
#   cold start, dengan dependensi   111.941  123.417        8583
\end{lstlisting}

\begin{table}[!htb]
  \centering
  \caption{Contoh hasil Latihan 5 yang sudah terisi}
  \label{tab:c14-lat5-contoh}
  \begin{tabularx}{\textwidth}{Xlll}
    \hline
    \textbf{Cara} & \textbf{p50 ms} & \textbf{p95 ms} & \textbf{Kali lipat} \\
    \hline
    \textit{warm} & 0,012 & 0,012 & 1 \\
    \hline
    \textit{cold start} & 19,069 & 22,248 & 1.610 \\
    \hline
    \textit{cold start}, 10 dependensi & 105,565 & 118,585 & 8.911 \\
    \hline
    \textit{cold start}, 13 dependensi & 111,941 & 123,417 & 8.583 \\
    \hline
  \end{tabularx}
\end{table}

\begin{table}[!htb]
  \centering
  \caption{Lembar isian hasil pengukuran Latihan 5}
  \label{tab:c14-lat5-hasil}
  \begin{tabularx}{\textwidth}{Xlll}
    \hline
    \textbf{Cara} & \textbf{p50 ms} & \textbf{p95 ms} & \textbf{Kali lipat} \\
    \hline
    \textit{warm} & \dots & \dots & \dots \\
    \hline
    \textit{cold start} & \dots & \dots & \dots \\
    \hline
    \textit{cold start}, 10 dependensi & \dots & \dots & \dots \\
    \hline
    \textit{cold start}, 13 dependensi & \dots & \dots & \dots \\
    \hline
  \end{tabularx}
\end{table}

Baris terakhir contoh memuat kejanggalan. Tiga dependensi tambahan menaikkan
p50 menjadi 111,941 milidetik, tetapi kolom kali lipatnya justru turun dari
8.911 menjadi 8.583. Kedua baris berasal dari dua kali jalan yang berbeda,
dan pembaginya, yaitu p50 \textit{warm}, naik dari 0,012 menjadi 0,013
milidetik pada jalan kedua. Penyebut yang sekecil itu membuat kolom kali
lipat tidak dapat dibandingkan antarjalan, sehingga yang dibandingkan adalah
kolom milidetiknya.
